\documentclass[handout]{beamer}
\usepackage{graphicx}
\usepackage{xcolor}
\usepackage{multirow}
\setbeamertemplate{blocks}[rounded][shadow=true]
\setbeamercolor{block body}{bg=gray!20, fg=black}
\setbeamercolor{block title}{bg=blue!20, fg=black}

\title{Clase 2: Estrategias de Decisión y Horizontes}
\subtitle{Métodos para la Gestión - Business School}
\author{Daniela Pucci de Farias}
\institute{Universidad Torcuato Di Tella}
\date{Junio 2025}

\begin{document}

\frame{\titlepage}

\begin{frame}{¿Por qué cambian las estrategias?}
\begin{itemize}
    \item \textbf{Caso de estudio:} Precios Dinámicos en Aerolíneas.
    \item ¿Por qué el precio de un pasaje sube drásticamente 7 días antes del vuelo?
    \item Factores en juego:
    \begin{itemize}
        \item Inventario (asientos restantes).
        \item Tiempo restante (horizonte $T$).
        \item Perfil del cliente (Turista vs. Corporativo).
    \end{itemize}
\end{itemize}
\end{frame}

\begin{frame}{Modelando la Estrategia: MDP}
\begin{itemize}
    \item Un Proceso Markoviano de Decisión (MDP) se define por $(S, A, P, r)$:
    \begin{itemize}
        \item \textbf{Estado ($S$):} Situación actual (ej. 50 asientos libres, 10 días para el vuelo).
        \item \textbf{Acción ($A$):} Decisión (ej. fijar precio en \$200).
        \item \textbf{Transición ($P$):} Probabilidad de vender un pasaje.
        \item \textbf{Recompensa ($r$):} Ganancia inmediata.
    \end{itemize}
\end{itemize}
\end{frame}

\begin{frame}{La Lógica de "Mirar hacia Adelante"}
\begin{itemize}
    \item La \textbf{Ecuación de Bellman} resume el sentido común empresarial:
    \begin{block}{Valor = Hoy + Mañana}
    $$V_t^*(x) = \max_{a \in A} \{ \text{Recompensa Inmediata} + \text{Valor Futuro Esperado} \}$$
    \end{block}
    \item El impacto del \textbf{Horizonte $T$}:
    \begin{itemize}
        \item En logística con fecha de entrega o campañas electorales, el tiempo dicta la agresividad de la acción.
    \end{itemize}
\end{itemize}
\end{frame}

\begin{frame}{Laboratorio: Perfiles de Cliente}
\begin{itemize}
    \item \textbf{Ejercicio computacional:}
    \item Modificar el código Python para dos tipos de pasajeros:
    \begin{enumerate}
        \item \textbf{Turista:} Alta sensibilidad al precio, compra con anticipación.
        \item \textbf{Profesional:} Baja sensibilidad al precio, compra a última hora.
    \end{enumerate}
    \item \textbf{Objetivo:} Observar cómo la política óptima $u^*(x,t)$ ajusta los precios automáticamente según la mezcla de clientes.
\end{itemize}
\end{frame}

\end{document}
